\title{LongRoPE: Extending LLM Context Window Beyond 2 Million Tokens}

\begin{abstract}
A substantial context window is an advantageous characteristic for large language models (LLMs). Nevertheless, existing methods for expanding context length are generally constrained to approximately 128k tokens. This limitation stems from expensive fine-tuning procedures, limited availability of lengthy text data, and the presence of problematic values arising from novel token positions.
\end{abstract}

This study presents LongRoPE, which, for the first time, expands the context window of pre-trained LLMs to a remarkable \textbf{2048k} tokens using at most 1,000 fine-tuning steps, with training lengths up to 256k, all while preserving the model's original performance on short contexts. The proposed approach hinges on three main contributions: (i) We uncover and harness two distinct non-uniform patterns in positional interpolation via an efficient search process, enabling improved initialization for fine-tuning and supporting an 8$\times$ context window extension without any fine-tuning; (ii) We propose a staged extension method whereby a 256k-length LLM is fine-tuned first, and then a subsequent round of positional interpolation is applied to this extended model to reach a 2048k token context window; (iii) We recalibrate LongRoPE on sequences of 8k length to fully restore short context window performance. Comprehensive evaluations conducted on LLaMA2 and Mistral, spanning multiple tasks, verify the efficacy of our technique. LongRoPE-extended models maintain the original neural architecture with only slight adjustments to the positional embeddings, thereby allowing compatibility with most existing optimizations. The code will be released at \texttt{https://github.com/microsoft/LongRoPE}

\section{Introduction}

Although Large Language Models (LLMs) have achieved impressive results across a range of tasks~\cite{gpt4,llama2}, they are constrained by a finite context window size—for instance, LLaMA2 is restricted to 4096 tokens~\cite{llama2}. When input exceeds this window, LLMs exhibit reduced effectiveness, largely due to their lack of exposure to such extended positional contexts during training. This limitation becomes particularly problematic in scenarios such as in-context learning with many provided examples~\cite{Huang2023FewerIM} or when deploying LLM agents~\cite{park2023generative,madaan2023selfrefine}.

Recent research has demonstrated that it is feasible to expand a pre-trained LLM’s context window up to approximately 128k tokens by performing fine-tuning with lengthier documents~\cite{longlora,pi,yarn,zhang2024soaring,liu2023scaling}. Nevertheless, scaling the context window further encounters three primary challenges. Firstly, incorporating novel, previously unseen position indices can introduce severe distributional anomalies, which significantly impedes the fine-tuning process and may prevent successful convergence~\cite{pi}. This problem becomes even more pronounced when scaling from 4k to over 1000k positions, since more than 90\% of the indices would be entirely new. Secondly, effective fine-tuning depends on the availability of appropriately long textual data. Yet, in existing datasets, texts surpassing 1000k tokens are rare, and training on such extended sequences incurs substantial computational costs, making it difficult to manage both in terms of time and required GPU capacity. Lastly, the shift to ultra-long context windows disperses the attention mechanism, distributing it across an expansive range of token positions, which ultimately undermines the model’s effectiveness on the original, shorter contexts~\cite{pi}.

A common strategy to address the initial issue is to perform interpolation on RoPE positional embeddings~\cite{rope,pi}, adjusting new positional indices to fit within the range seen during pretraining, as illustrated in Fig.\ref{fig:overview}. Position Interpolation (PI)~\cite{pi} modifies RoPE’s rotary angles using a linear scaling based on the extension ratio. The NTK method~\cite{ntk,dynamicntk} employs non-uniform interpolation and extrapolation across the dimensions of RoPE. YaRN~\cite{yarn} divides RoPE’s dimensions into three groups according to frequency, applying extrapolation, NTK, and linear interpolation techniques to each group respectively. Nonetheless, within the Transformer architecture, positional embeddings possess \emph{complex, non-uniform information entropy}. Current methods do not adequately exploit this subtle heterogeneity, which results in some information being lost and restricts the attainable length of the context window.

Section~\ref{sec:analysis} demonstrates empirically two principal insights: \textbf{(1)} Successful positional interpolation requires accounting for two types of non-uniformity: differences across RoPE dimensions and across token positions. Reduced interpolation benefits both lower RoPE dimensions and early token positions; however, the optimal strategy is contingent upon the specific desired extension length. \textbf{(2)} By integrating these non-uniformities into the positional interpolation process, it becomes possible to preserve vital aspects of the original RoPE encoding, especially regarding essential dimensions and positions of tokens. This approach reduces the degradation of information caused by interpolation and therefore offers a stronger initialization for subsequent fine-tuning. Furthermore, it enables models to be extended up to 8$\times$ longer even without additional fine-tuning.

Building on these insights, we propose \textbf{LongRoPE}, a robust approach that expands the context window of LLMs to accommodate over 2 \emph{million} tokens. The design of LongRoPE centers on three principal innovations. To begin with, {LongRoPE} leverages multidimensional non-uniformities present in positional interpolation. Specifically, it determines optimal rescale factors for the rotation angles within each dimension of RoPE by analyzing token positions. Since the search space for these rescale factors increases exponentially relative to the desired extension ratio, {LongRoPE} employs an evolutionary search algorithm enhanced by two optimization strategies, thereby significantly improving search efficiency. Figure~\ref{fig:overview} illustrates a representative example of the rescaled RoPE discovered through this process.

Subsequently, {LongRoPE} utilizes a resource-efficient and incremental extension method to reach a 2048k context window, circumventing the necessity for explicit fine-tuning on ultralong documents, which are seldom available. The approach initiates by identifying a suitable 256k context length within the pre-trained LLM, followed by fine-tuning the model at this window size. Given that our non-uniform positional interpolation technique supports up to an 8$\times$ increase in window length without fine-tuning, we proceed to a secondary phase that involves searching for updated RoPE rescaling factors on this previously fine-tuned, length-extended LLM. This two-stage process enables the realization of a 2048k context window for both LLaMA2 and Mistral~\cite{mistral}.

Lastly, to prevent loss of performance on the model's native (shorter) context window, {LongRoPE} further tunes the RoPE rescaling coefficients within the extended LLM. In an approach analogous to upscaling from 256k to 2048k, our method also performs a downscaling process to 4k and 8k context lengths on the LLM that was fine-tuned at 256k, utilizing a search algorithm to reduce the degree of positional interpolation. During inference, whenever the input sequence is under 8k tokens in length, RoPE is adjusted using the optimized rescale factors obtained from our search process.

A comprehensive set of experiments involving multiple LLM architectures and a range of long-context tasks highlights the strong performance of our approach. Our results indicate that LongRoPE consistently preserves low perplexity across evaluation lengths spanning from 4k up to 2048k, surpasses 90\% accuracy in passkey retrieval, and achieves accuracy on standard benchmarks comparable to those crafted for a 4096-token context window. The method is broadly applicable to all LLMs utilizing RoPE embeddings. We intend to make both our implementation and the LongRoPE-2048k models publicly available.

\section{Non-uniformity in Positional Interpolation}
\label{sec:analysis}

\subsection{Preliminary}

Transformer architectures necessitate the incorporation of explicit positional data—typically through position embeddings—to encode the sequence order of input tokens. This study centers on the RoPE~\cite{rope} position embedding mechanism, which has been broadly adopted in state-of-the-art LLMs. Specifically, the RoPE encoding assigned to a token at position index $n$ can be succinctly represented as:

\begin{equation}
		\fontsize{7.8}{7.8} \selectfont
	\label{eq:rope}
	\left[ cos(n\theta_0), sin(n\theta_0), cos(n\theta_1), \cdot\cdot\cdot, cos(n\theta_{d/2-1}), sin(n\theta_{d/2-1}) \right]   
\end{equation}

where $d$ denotes the embedding dimensionality, 
$n\theta_i$ refers to the rotational angle assigned to the token located at position $n$, and 
$\theta_i=\theta^{-2i/d}$ specifies the rotation frequencies. Within RoPE, $\theta$ is typically set to 10000 by default.

\textbf{\emph{Context window extension ratio $s$} and positional interpolation}. We denote $s$ as the quotient of the extended context length $L'$ over the original context length $L$, i.e., $s=\frac{L'}{L}$.

When increasing the context window from $L$ to $L'$, prevailing positional interpolation techniques propose modifying the rotation frequencies $\theta_i$ by scaling them according to the extension ratio $s$. Defining $\beta=\theta^{2/d}$, and letting $\lambda$ represent the effective rescaling factor corresponding to $s$, we can generalize these various positional interpolation methods with the following formulation:

\begin{equation}	
	\fontsize{7.5}{7.5} \selectfont
	\label{eq:unified_rope}
	\begin{aligned}
		\left[cos\left(\frac{n}{\lambda(\beta)^0}\right), sin \left(\frac{n}{\lambda(\beta)^0}\right), cos\left(\frac{n}{\lambda(\beta)^1}\right), \cdot\cdot\cdot,  sin \left(\frac{n}{\lambda(\beta)^{d/2-1}}\right) \right]   
	\end{aligned}
\end{equation}

\textbf{Linear positional interpolation (PI)}. PI~\cite{pi} introduces a method where positional indices are linearly interpolated as long as they remain within the original sequence length used during pre-training. When extending to a sequence scaled by a factor $s$, all rotation angles corresponding to position encodings are uniformly scaled down by $\lambda = s$ for each RoPE dimension. This scaling results in positional encodings being packed too closely together, which can impair the model’s capability to differentiate between tokens with similar positions. Consequently, the effectiveness of PI diminishes, particularly at larger extension ratios.

\textbf{NTK-based interpolation and extrapolation}. ~\cite{ntk,dynamicntk} approach RoPE from the standpoint of information encoding, leveraging Neural Tangent Kernel (NTK) theory~\cite{jacot2018neural,tancik2020fourier}. To address the issue of positional crowding in PI, these works propose redistributing interpolation stress across the different RoPE dimensions: lower (higher frequency) dimensions are scaled to a lesser extent, while higher (lower frequency) dimensions receive greater scaling. This approach enables both positional interpolation and extrapolation, with $\lambda=s^{i}$. The dynamic NTK approach~\cite{dynamicntk} further refines this by dynamically modulating the extension ratio at each position according to the current sequence length. In contrast to PI, which typically requires model fine-tuning, NTK-informed strategies permit expansion of the context window in settings where fine-tuning is not performed, albeit typically up to a 4$\times$ extension limit.

\textbf{YaRN}~\cite{yarn} brings notable enhancements to the effectiveness of positional interpolation. The approach segments the RoPE dimensions into three separate groups according to frequency, with each group subjected to a distinct interpolation method. Dimensions associated with high frequencies are extrapolated ($\lambda$=1), while those corresponding to low frequencies are linearly interpolated (PI). For RoPE dimensions that are intermediate in frequency, the NTK strategy is applied. Central to YaRN's methodology is the classification of RoPE dimensions into these groups, a process that is currently determined through manual, empirical tuning. As a consequence, this may not yield optimal outcomes for novel LLM architectures.

\subsection{Study on Non-uniform Positional Interpolation}

Drawing from the approaches of NTK and YaRN, it is evident that their performance improvements stem largely from incorporating non-linear components, notably through addressing varying frequency patterns within the RoPE dimensions for tailored interpolation and extrapolation. Nevertheless, existing methods for introducing non-linearity are predominantly governed by manually crafted heuristics. This prompts two primary inquiries: (1) Does present positional interpolation represent the best possible solution? (2) Might there exist alternative non-linear functions yet to be investigated?

In order to address these questions, we employ evolution search (refer to Sec.~\ref{sec:method}) to identify improved non-uniform positional interpolations tailored for LLaMA2-7B. The process is directed by evaluating perplexity, with five randomly chosen samples from the PG19~\cite{pg19} validation dataset. Our experimental investigation uncovers several important insights.

\textbf{Finding 1}: \textit{RoPE dimensions exhibit substantial non-uniformities, which are not effectively handled by current positional interpolation methods.}

We optimize $\lambda$ individually for each RoPE dimension as defined in Eq.~\ref{eq:unified_rope}. Table~\ref{tbl:exp1} presents a comparison of LLaMA2-7B’s perplexity, evaluated on PG19 and Proof-pile~\cite{proof-pile} test sets, across various interpolation methods in a zero-shot (non-fine-tuned) setting. The results indicate that our optimized configuration achieves marked gains in perplexity, implying that both current linear (PI) and alternative non-uniform interpolation approaches (Dynamic-NTK and YaRN) are not optimal. Interestingly, YaRN performs worse than both PI and NTK on the PG19 dataset, which is likely attributable to its failure to attain the full target context window in the absence of fine-tuning. Specifically, in an 8k context window, the perplexity for YaRN surges notably after reaching a length of 7k.

In our investigation, we observed that the rescaling factors $\lambda$ in Eq.~\ref{eq:unified_rope} are not uniform, which stands in contrast to the constant scaling factor $s$ used in PI, the formula in NTK, and the group-based computation in YaRN. The introduction of these variable scaling factors leads to substantial enhancements in LLaMA2's language modeling capabilities (measured by perplexity) at extended context lengths of 8k and 16k, all achieved without additional fine-tuning. This improvement is attributable to the positional embeddings' ability to better retain characteristics of the original RoPE, particularly in critical dimensions, thereby easing the challenge for LLMs when differentiating between nearby token positions.

\textbf{Finding 2}: \textit{RoPE for the initial tokens in the input sequence should be extrapolated with less interpolation.} 

We propose that the initial $\hat{n}$ tokens within input sequences should undergo minimal interpolation in their RoPE representations. This rationale stems from their characteristically high attention weights, emphasizing their significance in attention computations, a phenomenon noted in Streaming LLM~\cite{streamingllm} and LM-Infinite~\cite{han2023lminfinite}. To examine this hypothesis, we expand the context length to 8k and 16k using both PI and NTK strategies, and for each setting, refrain from applying interpolation to the first $\hat n$ tokens, varying $\hat n$ across the set (0, 2, ..., 256). When $\hat n$ equals 0, the setup defaults to standard PI or NTK. The results, summarized in Table~\ref{tbl:tokenposition}, yield two principal findings: \textbf{(1)} preserving position information of the initial tokens without interpolation leads to clear performance gains, and \textbf{(2)} the most effective value of $\hat n$ is contingent upon the desired length of context extension.

\textbf{Finding 3}: \textit{Non-uniform positional interpolation effectively extends LLM context window in both fine-tuning and non-fine-tuning settings.} 

Although our non-uniform position interpolation discovered via search notably enhances extension capability at 8k and 16k context lengths without additional training, extending to longer contexts necessitates fine-tuning. Therefore, we apply fine-tuning to LLaMA2-7B equipped with our searched RoPE configuration for a 64k context window (details provided in the Appendix). As indicated in Table~\ref{tbl:finetune}, our approach achieves substantial gains over both PI and YaRN, irrespective of whether LLaMA2-7B is evaluated pre- or post-fine-tuning. These improvements can be attributed to the robust use of non-uniform positional interpolation, which reduces information degradation and yields a superior starting point for fine-tuning.

\textbf{Summary}. In this work, we identify two forms of non-uniformity: variability across RoPE dimensions and across token positions. Strategic exploitation of these non-uniformities within positional interpolation produces considerable enhancements in extending LLM context lengths.

\section{LongRoPE}

\label{sec:method}
Building on these observations, we propose LongRoPE, which initially incorporates an optimized search algorithm designed to leverage the two identified non-uniformities. Subsequently, this approach is employed to expand the context window of LLMs to exceed 2 million tokens.

\subsection{Problem Formulation}

The presence of these two forms of non-uniformity results in an extensive solution space and adds layers of complexity to the optimization process. To manage this, we recast the multidimensional non-uniform position interpolation optimization challenge into a search framework.

Given an LLM designed for a context window size of $L'$ and a set of lengthy input documents $\mathbf{X}$, where each sequence $\mathbf{x}\in\mathbf{X}$ contains more tokens than $L'$, let the original rotary angle in RoPE embeddings for the $i^{th}$ dimension at token position $n$ be represented as $\frac{n}{\beta^i}$. The corresponding optimization task can thus be stated as:

\begin{equation}
\begin{gathered}
		\label{eq:problem}

\underset {\mathbf{x}\in\mathbf{X}; \, |\mathbf{x}|\ge L'}{ \operatorname {arg\,min} } \,  \mathcal{L}\, \left( \text{LLM} (\text{RoPE}, \mathbf{X})\right),	\text{where \;} 
\\
\scalebox{0.93}{
$\underset {\begin{subarray}{l} i=0,\ldots,\frac{d}{2}-1;\\ n\in[ 0,|\mathbf{x}|);\end{subarray}}{\text{RoPE}_(n)}= {\left[\ldots, \cos\left(\mathbb{I}(\hat{\lambda}_{i}, \hat{n})\frac{n}{\beta^i}\right), \sin\left(\mathbb{I}(\hat{\lambda}_{i}, \hat{n})\frac{n}{\beta^i}\right),\ldots\right] }$}\\
	\text{where \;} \mathbb{I}(\hat{\lambda}_{i},\hat{n})=\begin{cases}
	1&\mbox{\;  $n<\hat{n}$}\\
{\frac{1}{\lambda}_{i}}&\mbox{\; $n\geq\hat{n}$}
\end{cases}
	\end{gathered}
\end{equation}

In this context, we define a collection of rescale factors, $\mathbb{I}(\hat{\lambda}_{i},\hat{n})$, to address both types of non-uniformity. Here, $\hat{\lambda_i}$ corresponds to the dimension-specific non-uniformity in RoPE, while $\hat{n}$ represents non-uniformity linked to token positions. The factor $\mathbb{I}(\hat{\lambda}_{i},\hat{n})$ is used to adjust the rotational angle in the $i^{th}$ dimension of RoPE. Specifically, $\hat\lambda_{i}$ acts as the scaling coefficient, and $\hat{n}$ determines a token position threshold. For token positions earlier than $\hat{n}$ (that is, the first $\hat{n}$-1 tokens), the scaling coefficient $\hat\lambda_{i}$ is not applied, and the standard RoPE rotation angle $\frac{n}{\beta^i}$ remains unchanged. Once tokens are at positions $n\ge\hat{n}$, the rescaling factor comes into effect.

Assuming a desired context window length of $L'$, our goal is to determine the best set of rescale factors ($\mathbb{I}(\hat{\lambda}_{0},\hat{n})$, $\mathbb{I}(\hat{\lambda}_{1},\hat{n})$, ... $\mathbb{I}(\hat{\lambda}_{i},\hat{n})$, ...) corresponding to each RoPE dimension from the $1^{st}$ up to the $d^{th}$. By applying these rescale factors to the target $\text{LLM}$'s $\text{RoPE}$, we aim for the model to achieve minimal next-token prediction loss $\mathcal{L}$—that is, reduced perplexity—when presented with input sequences $\mathbf{X}$ of length $L'$.

\subsection{Searching the Non-uniform Position Interpolation}

To address the challenge formulated in Eq.~\ref{eq:problem}, we present an efficient and straightforward technique aimed at determining the optimal RoPE rescale factors. This approach is tailored to leverage the multidimensional, non-uniform characteristics inherent to position embeddings.

\textbf{Search space}. Our method encompasses a comprehensive search space that incorporates both sources of non-uniformity, as detailed in Table~\ref{tbl:searchspace}. In particular, the method permits the individual optimization of rescale factors for each dimension within the RoPE embedding. For the sake of simplifying the design of the search space, the algorithm searches over $\lambda_i$ and $\hat{n}$ in place of directly optimizing $\mathbb{I}(\hat{\lambda}_{i},\hat{n})$, with $\hat{\lambda}_{i}=1/\lambda_i$. As depicted in Table~\ref{tbl:searchspace}, each $\lambda_i$ is allowed to range from a minimum of 1.0 (corresponding to pure extrapolation) up to a maximum of $s\times1.25$ (representing a broader interpolation than PI), using a step increment of 0.01, where $s$ denotes the desired context window extension ratio.

The parameter $\hat{n}$ specifies how many of the earliest token positions preserve the original RoPE embedding, foregoing position interpolation. In practice, we select $\hat{n}$ from the set \{0, 1, 2, 4, 8, 12, 16, 20, 24, 28, 32, 64, 128, 256\} during the search process. If $\hat{n}=0$, then the learned rescale factors are applied across every token position.

\textbf{Evolution-based search}. The search space detailed in Table~\ref{tbl:searchspace} encompasses a vast array of positional interpolation configurations, making exhaustive exploration computationally prohibitive. For instance, selecting a $s=4\times$ extension yields $400^{128/2}\times14$ or approximately 4$\times10^{167}$ possible combinations. As the extension ratio increases, the number of potential solutions rises exponentially. To manage this complexity, we adopt evolution search~\cite{guo2020single} and integrate two optimization strategies that substantially enhance the efficiency of the search process. The comprehensive procedure is outlined in Algorithm~\ref{alg:search}.

\textit{Optimized initial population generation}. Rather than starting with an entirely random set of $P$ rescale factors, our approach incorporates the RoPE rescale factors associated with PI, NTK, and YaRN directly as three members of the initial population. The other $P$-3 individuals are produced by applying random mutations to these three rescale factors, each with a probability $p$.

\textit{Monotonically non-decreasing constraint}. Once the initial population has been produced, we evaluate the LLM perplexity for every individual. This involves applying each individual's specific RoPE rescale factors to the target LLM and then measuring the perplexity for input $\mathbf{X}$. The highest-performing $k$ individuals are selected as parents to proceed in the evolutionary process. Nevertheless, the extensive size of the search space means that straightforward mutation and crossover may frequently generate low-quality candidates, which leads to redundant perplexity evaluations. This inefficiency is exacerbated as $L'$ increases, since each perplexity assessment is computationally intensive.

To tackle this issue, we enforce a monotonic non-decreasing condition on the rescaling factors used for RoPE: $\lambda_i\le \lambda_{i+1}$. Only those RoPE configurations that comply with this monotonicity constraint are utilized in the perplexity assessments of LLMs, thereby streamlining the search process. Concretely, we stipulate that the values of $\lambda_i$ must increase monotonically as the RoPE dimension index $i$ progresses from $0$ to $63$. This requirement for dimensional monotonicity draws on the principles of NTK theory~\cite{jacot2018neural,tancik2020fourier,ntk}, which indicates that dimensions corresponding to higher frequencies (lower $i$) need less interpolation (i.e., smaller $\lambda_i$), whereas dimensions with lower frequencies (higher $i$) are suitable for greater interpolation (i.e., larger $\lambda_i$).

\begin{algorithm}[t]
	\small
	\caption{The search algorithm}
	\textbf{Input:} target LLM, input samples $\mathbf{X}$, population size $P$, mutation size $N_1$, crossover size $N_2$, maximum iterations $\mathcal{T}$, mutation probability $p$ \\

	\begin{algorithmic}[1]
		\label{alg:search}
		\STATE Top-k $\gets \phi$;
		\STATE $\text{P}_0 \gets$ \textit{Initialize\_population\_with\_optimization} ($P$, $\mathbf{X}$, $p$);
		\FOR{$i$ from $1$ to $\mathcal{T}$}
		\STATE \textit{Compute\_perplexity}(LLM, $\text{P}_{i-1}$, $\mathbf{X}$);
		\STATE Top-k $\gets$ \textit{Update\_Topk}(Top-k, $\text{P}_{i-1}$);
		\STATE $\text{P}_{mutation} \gets$ \textit{Mutation\_with\_mono\_constraint}(Top-k, $N_1$, $p$);
		\STATE $\text{P}_{crossover} \gets$ \textit{Crossover\_with\_mono\_constraint}(Top-k, $N_2$);
		\STATE $\text{P}_i = \text{P}_{mutation} \cup \text{P}_{crossover} \cup$ Top-k;
		\ENDFOR
		\STATE Output the member of Top-k exhibiting the minimum perplexity;
	\end{algorithmic}
\end{algorithm}

\textbf{8$\times$ extension without fine-tuning}. Through evolutionary search, our approach successfully determines optimal non-uniform RoPE rescale factors, selectively maintaining important dimensions and positions to reduce information loss caused by interpolation. As illustrated in Fig.\ref{fig:non-ft}, this technique allows us to expand LLaMA2's context window from 4k to 32k tokens without the need for fine-tuning. In comparison, prior approaches such as PI, and non-uniform NTK and YaRN, result in a significant increase in perplexity once the context window is extended beyond 2$\times$.

\subsection{Extending LLM Context Window to 2048K}
\label{sec:extendto2048k}

\textbf{Progressive extension to 2048k}. In this section, we present our strategy for scaling the context window of pre-trained LLMs beyond the conventional 4k limit to more than 2048k tokens. Our non-uniform positional interpolation approach enables up to an 8$\times$ increase in context length without the need for additional fine-tuning. However, for more substantial expansions (such as 512$\times$), fine-tuning becomes indispensable. A common approach involves identifying suitable RoPE rescaling parameters for the targeted 2048k window, followed by model fine-tuning. Nonetheless, this method encounters significant obstacles, notably the extremely high computational cost associated with training at such scales. Furthermore, our practical observations suggest that effectively fine-tuning LLMs for such large extension ratios presents considerable difficulty (refer to Appendix for further discussion).

Luckily, LongRoPE proves successful when applied to both base models and those that have undergone extension via fine-tuning. To this end, we propose an efficient and incremental approach, reaching the desired 2048k context window in only 1k fine-tuning iterations, all constrained to a maximum training sequence length of 256k.

$\Diamond$ \textit{LongRoPE search for scaling pre-trained LLMs to 256k.} Using LLaMA2 as a representative case, we perform a search targeting context window lengths of 128k and 256k. At this phase, the extension ratios achieved are 32$\times$ and 64$\times$, accordingly.

$\Diamond$ \textit{Fine-tuning to 256k}. Following pre-training, we further adapt the LLM to support a 256k context window through fine-tuning. The procedure involves initially fine-tuning LLaMA2 for 400 steps with the RoPE rescaled factors set for 128k. Upon completing these steps, we switch the RoPE rescaled factors to those corresponding to 256k and continue with 600 more fine-tuning steps from this checkpoint. Compared to a direct fine-tuning to 256k, this staged approach demonstrates greater efficiency.

$\Diamond$ \textit{Expanding the context length of a fine-tuned extended LLM to 2048k using LongRoPE search}. In the final step, we apply an additional search process to the LLM that has already been fine-tuned on 256k-length sequences. This approach enables the model to reach a significantly larger context window of 2048k tokens, all without requiring additional fine-tuning procedures. As a result, the model achieves an overall context extension of $512\times$.

\textbf{Degradation in shorter context window performance}. When expanding to a substantially longer 2048k context window, we observe that model performance declines within the initial context range. This decline is a recognized consequence of positional interpolation~\cite{pi}: the embedding for positions in the original context must now occupy a considerably smaller subspace of the position embedding, especially at higher dimensions. Consequently, this constriction impairs the model’s effectiveness. Specifically, at an extension ratio of 512$\times$, positions inside the original 4k context window are densely packed together.

To address this issue, we conduct an additional evolution search on the extended LLM, specifically tuning the RoPE rescale factors for shorter context lengths such as 4k and 8k. For these lengths, we lower the upper bound on the searched $\lambda$, as they necessitate less positional interpolation. When running inference, the LLM adapts the relevant RoPE rescale factors in real time.

\section{LongRoPE}

\label{sec:method}
Building on these insights, we propose LongRoPE, which begins by incorporating an effective search algorithm designed to leverage both non-uniformities to their fullest. Subsequently, this approach is employed to expand the context window of LLMs to exceed 2 million tokens.

\subsection{Problem Formulation}

These two types of non-uniformity significantly expand the solution space, thereby complicating the optimization process. To tackle this challenge, we recast the multidimensional non-uniform position interpolation optimization as a search problem.

Consider a LLM designed for a context window of size $L'$ and input documents $\mathbf{X}$ in which each segment $\mathbf{x}\in\mathbf{X}$ exceeds $L'$ tokens in length. For such cases, let the standard rotary angle for the $i^{th}$ dimension in the RoPE embedding at token position $n$ be expressed as $\frac{n}{\beta^i}$. The optimization objective is accordingly specified as follows:

\begin{equation}
\begin{gathered}
		\label{eq:problem}

\underset {\mathbf{x}\in\mathbf{X};\,|\mathbf{x}|\ge L'}{ \operatorname {arg\,min} }\,\mathcal{L}\left( \text{LLM} (\text{RoPE}, \mathbf{X})\right), \text{where}
\\
\scalebox{0.93}{
$\underset {\begin{subarray}{l} i=0, \ldots, \frac{d}{2}-1; \\ n\in[0,|\mathbf{x}|); \end{subarray}}{\text{RoPE}_{(n)}}={\left[ \ldots, \cos\left( \mathbb{I}(\hat{\lambda}_{i}, \hat{n}) \cdot \frac{n}{\beta^i} \right),\ \sin\left( \mathbb{I}(\hat{\lambda}_{i}, \hat{n}) \cdot \frac{n}{\beta^i} \right),\ldots \right]}$}
\\
\text{where} \quad \mathbb{I}(\hat{\lambda}_{i}, \hat{n}) = \begin{cases}
	1 &\text{if } n < \hat{n} \\
	\frac{1}{\lambda_{i}} &\text{if } n \geq \hat{n}
\end{cases}
\end{gathered}
\end{equation}

In this approach, we define a group of scaling factors, $\mathbb{I}(\hat{\lambda}_{i},\hat{n})$, to address the two types of non-uniformities present. Here, $\hat{\lambda}_{i}$ refers to the scaling of RoPE dimensions, while $\hat{n}$ pertains to non-uniformities across token positions. The term $\mathbb{I}(\hat{\lambda}_{i},\hat{n})$ is used to adjust the rotation angle for the $i^{th}$ dimension in the RoPE embedding, where $\hat{\lambda}_{i}$ represents the scaling parameter and $\hat{n}$ serves as the positional threshold. For the first $\hat{n}$-1 token positions, the scaling $\hat{\lambda}_{i}$ is not applied, and the standard RoPE rotary angle $\frac{n}{\beta^i}$ is maintained. Once the position $n$ reaches or exceeds $\hat{n}$, the scaling factor is introduced.

Assuming a desired context window length of $L'$, the task is to determine the best set of rescale factors ($\mathbb{I}(\hat{\lambda}_{0},\hat{n})$, $\mathbb{I}(\hat{\lambda}_{1},\hat{n})$, ... ,$\mathbb{I}(\hat{\lambda}_{i},\hat{n})$, ...) corresponding to the $1^{st}$ through $d^{th}$ dimensions of RoPE. By appropriately rescaling $\text{RoPE}$ in the target $\text{LLM}$, our goal is to reduce the next token prediction loss, $\mathcal{L}$ (such as perplexity), on input sequences $\mathbf{X}$ of length $L'$ as much as possible.

\subsection{Searching the Non-uniform Position Interpolation}

To address the challenge presented in Eq.~\ref{eq:problem}, we propose a straightforward but powerful approach that identifies optimal RoPE rescale factors, allowing us to leverage multidimensional and non-uniform characteristics in positional embeddings to their fullest extent.

\textbf{Search space}. Our method constructs an expansive search space designed to accommodate both forms of non-uniformity. As summarized in Table~\ref{tbl:searchspace}, this space permits the assignment of distinct rescale factors to each RoPE embedding dimension. To streamline the search space, our strategy focuses on varying $\lambda_i$ and $\hat{n}$, foregoing a direct search for $\mathbb{I}(\hat{\lambda}_{i},\hat{n})$, where $\hat{\lambda}_{i}=1/\lambda_i$. According to Table~\ref{tbl:searchspace}, the permissible values for $\lambda_i$ range from a minimum of 1.0 (corresponding to straightforward extrapolation) up to a maximum of $s\times1.25$ (enabling more substantial interpolation than that used in PI), incremented by 0.01. Here, $s$ denotes the ratio by which the context window is to be extended.

The parameter $\hat{n}$ specifies how many initial token positions retain the unaltered RoPE embedding, foregoing position interpolation. In practice, we permit $\hat{n}$ to take values from the set \{0, 1, 2, 4, 8, 12, 16, 20, 24, 28, 32, 64, 128, 256\}. Setting $\hat{n}=0$ indicates that every token position applies the determined rescale factors.

\textbf{Evolution-based search}. The search space described in Table~\ref{tbl:searchspace} encompasses a vast array of positional interpolation options, making it highly challenging to navigate efficiently. For instance, when the extension ratio is $s=4\times$, the number of possible configurations reaches $400^{128/2}\times14=4\times10^{167}$. As the extension ratio increases, the size of the search space grows at an exponential rate. To cope with this complexity, we adopt evolution search~\cite{guo2020single} and incorporate two optimization strategies that substantially enhance the efficiency of the search. The full process is outlined in Algorithm~\ref{alg:search}.

\textit{Optimized initial population generation}. Rather than populating the initial pool of $P$ rescale factors solely at random, we explicitly include the three RoPE rescale factors associated with PI, NTK, and YaRN as members of the initial group. The other $P$-3 individuals are then generated by applying random mutations to these three rescale factors, each with probability $p$.

\textit{Monotonically non-decreasing constraint}. Once the initial population is formed, we evaluate the LLM perplexity of every candidate. This involves applying their respective RoPE rescale factors within the target LLM and then measuring perplexity on the input $\mathbf{X}$. The highest-performing $k$ candidates are selected as parents for the subsequent evolutionary phase. Due to the extensive search space, straightforward mutation and crossover operations may frequently yield suboptimal candidates, resulting in numerous superfluous perplexity evaluations. This inefficiency is exacerbated for larger values of $L'$, where each perplexity assessment requires significant computational resources.

To tackle this issue, we enforce a monotonic non-decreasing constraint on the chosen RoPE rescale factors such that $\lambda_i \le \lambda_{i+1}$. Only RoPE configurations adhering to this monotonicity condition are utilized for LLM perplexity assessment, which substantially cuts down the search space. More precisely, we stipulate that $\lambda_i$ must increase monotonically with respect to the RoPE dimension, where $i$ ranges from $0$ to $63$. The rationale for this monotonicity across dimensions draws on the NTK theory~\cite{jacot2018neural,tancik2020fourier,ntk}, which posits that lower-indexed dimensions—corresponding to higher frequencies—require less interpolation (i.e., smaller $\lambda_i$ values), whereas higher-indexed, lower-frequency dimensions benefit from greater interpolation (i.e., larger $\lambda_i$ values).

\begin{algorithm}[t]
	\small
	\caption{The search algorithm}
	\textbf{Input:} target LLM, input samples $\mathbf{X}$, population size $P$, mutation size $N_1$, crossover size $N_2$, maximum iterations $\mathcal{T}$, mutation probability $p$\\

	\begin{algorithmic}[1]
		\label{alg:search}
		\STATE Top-k $\leftarrow \phi$;
		\STATE $\text{P}_0$ $\leftarrow$ \textit{Initialize\_population\_with\_optimization} ($P$, $\mathbf{X}$, $p$);
		\FOR{$i = 1$ \textbf{to} $\mathcal{T}$}
		\STATE \textit{Compute\_perplexity} (LLM, $\text{P}_{i-1}$, $\mathbf{X}$);
		\STATE Top-k $\leftarrow$ \textit{Update\_Topk} (Top-k, $\text{P}_{i-1}$);
		\STATE $\text{P}_{mutation}$ $\leftarrow$ \textit{Mutation\_with\_mono\_constraint} (Top-k, $N_1$, $p$);
		\STATE $\text{P}_{crossover}$ $\leftarrow$ \textit{Crossover\_with\_mono\_constraint} (Top-k, $N_2$);
		\STATE $\text{P}_i$ $\leftarrow$ $\text{P}_{mutation}$ $\cup$ $\text{P}_{crossover}$ $\cup$ Top-k;
		\ENDFOR
		\STATE Return the member in Top-k exhibiting the lowest perplexity;
	\end{algorithmic}
\end{algorithm}

\textbf{8$\times$ extension without fine-tuning}. By leveraging evolutionary search, our approach efficiently discovers optimal non-uniform RoPE rescale factors, safeguarding crucial positions and dimensions to reduce information degradation from interpolation. As illustrated in Fig.\ref{fig:non-ft}, this technique enables LLaMA2's context window to be expanded from 4k up to 32k tokens without the need for fine-tuning. In comparison, current approaches like PI, as well as non-uniform NTK and YaRN, exhibit a significant increase in perplexity when attempting extensions beyond 2$\times$.

\subsection{Extending LLM Context Window to 2048K}
\label{sec:extendto2048k}

\textbf{Progressive extension to 2048k}. In this section, we present our approach for expanding the context window of pre-trained LLMs from the standard 4k to beyond 2048k tokens. Our results indicate that the use of non-uniform positional interpolation enables an 8$\times$ increase in context length without the need for fine-tuning. However, achieving more substantial extensions (such as 512$\times$) necessitates fine-tuning. A possible strategy involves determining suitable RoPE rescaling factors corresponding to the target 2048k context size prior to initiating fine-tuning. Nevertheless, this procedure is hindered by the substantial computational resources required. In addition, our observations suggest that fine-tuning LLMs at very large extension ratios presents considerable difficulties (refer to Appendix for details).

Fortunately, LongRoPE demonstrates efficacy in both pre-trained and fine-tuned extended LLM scenarios. Consequently, we propose a streamlined, incremental approach that attains the desired 2048k target using only 1k fine-tuning steps, each conducted at a maximum training length of 256k.

$\Diamond$ \textit{LongRoPE search for scaling pre-trained LLMs to 256k}. Using LLaMA2 for illustration, we perform a search procedure with intended context window lengths of 128k and 256k tokens. At this phase, the context is extended by factors of 32$\times$ and 64$\times$ respectively.

$\Diamond$ \textit{Fine-tuning to 256k}. Subsequently, the pre-trained LLM undergoes fine-tuning to support a context window length of 256k. The process begins by fine-tuning LLaMA2 for 400 steps utilizing the RoPE rescaled factors set for 128k. Following this stage, the RoPE rescaled factors are switched to those for 256k on the resultant checkpoint, after which fine-tuning continues for another 600 steps. This two-phase strategy is found to be more effective than attempting direct fine-tuning at the 256k context length.

$\Diamond$ \textit{LongRoPE search applied to extend fine-tuned LLMs to 2048k}. In the final stage, we apply a subsequent search process to the LLM that has already been fine-tuned for sequences of length 256k. Through this approach, we successfully achieve an exceptionally large context window of 2048k tokens, all without additional rounds of fine-tuning. The resulting expansion factor is $512\times$.

\textbf{Shorter context window recovery}.  When scaling up to a notably large 2048k context window, we observe a decline in model performance within the standard context window. This drawback has been previously documented in positional interpolation methods~\cite{pi}, where the embedding for positions in the original context range is compressed into a significantly smaller space, impairing model effectiveness. Specifically, at an extension ratio of 512$\times$, the positional encodings corresponding to the initial 4k context window are densely packed.

To address this issue, we conduct an additional evolution search on the expanded LLM, specifically tuning the RoPE rescale factors for relatively short context lengths (such as 4k and 8k). Since these lengths require less positional interpolation, we constrain the upper limit for the searched $\lambda$. When running inference, the LLM automatically adapts the appropriate RoPE rescale factors in real time.

 \section{Experiments}

\label{sec:exp}
\subsection{Setup}
\textbf{Evaluation Tasks and Models}. LongRoPE is integrated into LLaMA2-7B and Mistral-7B, and its effectiveness is assessed across three dimensions: (1) perplexity scores of the extended-context language models when processing lengthy documents; (2) performance on the Passkey retrieval task, designed to test the model’s proficiency in extracting a simple passkey embedded among a large amount of distractor text; and (3) accuracy on standard large language model benchmarks using a context window limited to 4096 tokens.

\textbf{Fine-tuning}. The LLaMA2 model is optimized using a learning rate of 2e-5 with a linear decay schedule, utilizing a global batch size of 32. Fine-tuning is performed for 400 steps on the Redpajama~\cite{redpajama} dataset, which is segmented into 128k-token blocks, each bounded by BOS and EOS tokens. Building upon the resulting checkpoint, we continue training for an additional 600 steps to reach a 256k-token context window. Training for the 128k context employs 8 A100 GPUs using the distributed framework~\cite{cube}, whereas extending to 256k context requires scaling to 16 A100 GPUs. For Mistral, the training procedure employs a fixed learning rate of 1e-6 with a global batch size of 64. Both the 128k and 256k context variants follow the YaRN~\cite{yarn} methodology, entailing 400 steps of fine-tuning on Together Computer’s Long-Data Collections~\cite{mistraldata} with a 16k sequence length, conducted using 4 A100 GPUs.

\textbf{Search}. When the target window size is up to 256k, we set the parameters as follows: $P$=64, $N_1$=$N_2$=16, $p$=0.3, $\mathcal{T}$=40. In each iteration, we select the top 32 candidates for mutation and crossover. Perplexity is computed on 5 randomly chosen samples from the PG19 validation dataset, ensuring each meets the minimum required length for the target context. For target windows greater than 512k, we reduce the sizes of the population, mutation, and crossover pools by half. In this case, perplexity evaluation uses 3 randomly selected examples from the Pile-Books3~\cite{pile} validation dataset.

\textbf{Baselines}. In order to achieve a context window of 2048k, we conducted fine-tuning of our models using window sizes of 128k and 256k tokens. As a result, we refer to the corresponding models as LongRoPE-2048k (ft=128k) for LLaMA2 and LongRoPE-2048k (ft=256k) for Mistral. These models are evaluated alongside leading baselines for context window extension, particularly those open-sourced LLMs that have undergone fine-tuning subsequent to positional interpolation through methods such as PI, NTK, and YaRN. The baseline models in our comparison set include Together-32k~\cite{together}, Code LLaMA~\cite{codellama}, LongLoRA-full-FT-100k~\cite{longlora}, YaRN-LLaMA, and YaRN-Mistral~\cite{yarn}.

\subsection{Main Results}

\label{sec:mainresult}
\textbf{Long sequence language modeling within 256k}. We first assess our approach in comparison to leading extended LLMs for sequence lengths up to 256k tokens. To highlight generalizability, our experiments utilize two datasets: the test partitions from Proof-pile~\cite{pg19} and PG19~\cite{pile}. Perplexity is measured across different context lengths, employing a sliding window of size 256. For PG19, evaluation is performed on the entire test set, consisting of 100 documents. In the case of Proof-pile, consistent with the protocol of YaRN~\cite{yarn}, we randomly choose 10 samples, each exceeding a minimum length of 128k.

Table~\ref{tbl:proof-pile} and Table~\ref{tbl:pg19} present a comparison of perplexity scores for LLaMA2 and Mistral, each augmented with various interpolation strategies, evaluated on Proof-pile and PG19 datasets, respectively. Two principal findings are emphasized: \textbf{(1)} the perplexity of our enhanced models generally declines as evaluation lengths increase from 4k to 256k, which demonstrates their capability to exploit longer sequences effectively. \textbf{(2)} Remarkably, despite operating with a context window that is 16$\times$ longer—a scenario that ordinarily impedes performance at shorter lengths—our LongRoPE-2048k models surpass current state-of-the-art baselines when evaluated at the 256k context length.

\textbf{Long sequence language modeling beyond 2000k}. To assess performance on exceptionally lengthy texts, we employ the Books3~\cite{pile} dataset. For the sake of efficient evaluation, a random subset of 20 books—each surpassing 2048k tokens—is selected, and a 256k sliding window is applied.

As illustrated in Table~\ref{tbl:books3}, LongRoPE effectively increases the context window of both LLaMA2-7B and Mistral-7B up to 2048k, while attaining perplexity results that match or surpass those of baseline methods in shorter windows ranging from 8k to 128k. Distinct differences in performance arise between LLaMA2 and Mistral at the 2048k context length. While Mistral achieves superior perplexity compared to baselines for shorter input lengths, its perplexity rises above 7 when the context exceeds 256k tokens. For LLaMA2, as anticipated, perplexity generally diminishes with increasing context window length, though it experiences slight upticks at 1024k and 2048k. Notably, for LLaMA2, fine-tuning LongRoPE-2048k at 256k yields better results than fine-tuning at 128k, attributed to the smaller secondary extension ratio (i.e., 8$\times$ versus 16$\times$). Conversely, for Mistral, fine-tuning at 128k leads to improved outcomes relative to 256k. This discrepancy mainly stems from the use of a 16k training length, consistent with YaRN's methodology, during both 128k and 256k fine-tuning of Mistral; this choice constrains Mistral's capacity for further extending the context window post-fine-tuning.

\textbf{Passkey retrieval}. Next, we examine how well models utilize context windows during generation by conducting the synthetic passkey retrieval task introduced by ~\cite{passkey}. In this benchmark, the model must extract a randomly generated five-digit passkey that is embedded somewhere within a lengthy document. Details of the prompt template can be found in the appendix. For our evaluation, we repeat the retrieval experiment 10 times, each with the passkey inserted at a uniformly random position throughout the chosen context length.

Fig.~\ref{fig:passkey} presents a comparative analysis of retrieval accuracy against baseline methods. The performance of existing models deteriorates sharply, with accuracy falling to zero once the sequence length surpasses 128k. Conversely, LongRoPE-LLaMA2-2048k (ft=256k) demonstrates robust retrieval capabilities, sustaining high accuracy rates ($\ge$90\%) consistently across input lengths ranging from 4k up to 2048k, even when tasked with extracting a passkey among millions of tokens. Similarly, LongRoPE-Mistral-2048k (ft=128k) achieves perfect accuracy up to 1800k tokens, though its performance declines to 60\% at 2048k tokens. This outcome is in agreement with the results in Table~\ref{tbl:books3}, which show a modest rise in perplexity at 2048k.

\textbf{Performance on established benchmarks with the original context length}. LongRoPE-2048k models are assessed on the original context window employing the Hugging Face Open LLM Leaderboard~\cite{huggingfaceboard}, across both zero-shot and few-shot configurations. For evaluation, we apply a 25-shot setup on ARC-Challenge~\cite{allenai:arc}, a 10-shot protocol for HellaSwag~\cite{zellers2019hellaswag}, a 5-shot format for MMLU~\cite{mmlu}, and a 0-shot scenario for TruthfulQA~\cite{lin2021truthfulqa}.

As illustrated in Table~\ref{tbl:commonsense}, our models deliver results on par with those observed on the initial benchmark, which was intended for a limited context window, and they surpass the original Mistral on TruthfulQA by an increase of +0.5\%. Although LongRoPE-LLaMA2-2048k—fine-tuned using a 256k context—exhibits a marginally higher drop in performance, its scores for the majority of tasks continue to fall within acceptable bounds.

\subsection{Ablation Results}

\textbf{Assessing the impact of a second positional interpolation}. Within our approach to progressive extension, we apply our search-based methodology to perform an additional round of non-uniform positional interpolation on LLMs that have already undergone fine-tuning. To evaluate the benefit of this procedure, we experiment with our fine-tuned LLaMA2-256k model, further extending its context length to 512k, 1024k, and 2048k using both PI and YaRN methods. According to the results summarized in Table~\ref{tbl:secondsearch}, our non-uniform positional interpolation preserves stable perplexity levels across larger extensions. Conversely, perplexity with PI and YaRN escalates rapidly as the extension ratio increases.

\textbf{Effectiveness of recovery at shorter context lengths.} In order to address performance degradation at reduced context lengths, we recalibrate the RoPE factors for LongRoPE-2048k using our optimization algorithm. In this process, we restrict the maximum permitted scale factors during the search, which leads to diminished interpolation when evaluating at shorter context lengths such as 4k and 8k. Table~\ref{tbl:shortperformance} presents a comparison of perplexity for LongRoPE-LLaMA2-2048k on the Proof-pile dataset at 4k and 8k contexts, as well as the corresponding mean LLM benchmark accuracy. The data illustrate that there is a marked increase in performance at these shorter context lengths.

\textbf{Analysis on the two forms of non-uniformities}. In this section, we conduct an ablation study on the two identified non-uniformities to determine their individual effects on model performance. The experimental design involves two scenarios: (i) scaling up LLaMA2-7B to relatively short sequence lengths (16k and 32k), where we employ three strategies—PI, modifying only the RoPE dimension, and modifying both non-uniformities; and (ii) increasing the context window of our fine-tuned 256k-length LLaMA2 to a substantially longer 2048k, applying the same variations. We report perplexity results prior to any additional fine-tuning. As presented in Table~\ref{tbl:searchdimension}, introducing non-uniformity to the RoPE dimension yields a notable reduction in perplexity relative to PI’s linear interpolation baseline. Altering non-uniformity in token positions leads to further performance improvements at 16k and 32k lengths, although this effect is not evident at 2048k, potentially due to the extraordinary sequence length. The approach of retaining only the earliest tokens without any interpolation appears ineffective under these conditions, suggesting an avenue for further investigation.

\section{Related Works}

Besides position interpolation-based techniques, this section also reviews alternative related methodologies.

\textbf{Retrieval-based} methods leverage external memory systems to store extensive contextual information from previous segments and incorporate retrieval modules to access pertinent documents during inference~\cite{longllama,wang2023augmenting,borgeaud2022improving}. Such strategies often require substantial changes to the underlying LLM architectures. By comparison, our approach introduces only minimal alterations to positional embeddings, offering a more streamlined solution. Furthermore, our method accommodates a wider range of long-context applications, including tasks like long document summarization and few-shot learning, extending beyond the typical scope of retrieval-focused systems.

\textbf{Attention-based context window extensions}. In addition to the use of positional embedding interpolation, alternative approaches extend the input context length within the standard LLM context window by adapting attention mechanisms~\cite{han2023lminfinite, streamingllm,parallelwindow}. The central concept involves employing innovative attention masks to address the attention explosion problem introduced by extended positional indices. Such strategies can be seen as complementary to methods based on positional interpolation.

\textbf{Fine-tuning based approaches} concentrate on optimizing the fine-tuning process of pre-trained LLMs by altering position embeddings to support longer context windows. For example, Code LLaMA~\cite{codellama}, LLaMA2 Long~\cite{xiong2023effective}, and ScaledRoPE~\cite{liu2023scaling} implement a significantly larger RoPE base value and then conduct fine-tuning at the intended sequence length. Our approach, in contrast, allows for adaptation across a wide range of target sequence lengths and can scale to lengths exceeding 2M tokens. Recently, as fine-tuning LLMs for extended contexts (exceeding 128k tokens) has become increasingly resource-intensive, solutions such as LongLoRA~\cite{longlora} and PoSE~\cite{zhu2023pose} have been developed to reduce computational demands. Our proposed method is complementary and can be combined with these efficient fine-tuning techniques.

\section{Conclusion}

This paper introduces LongRoPE, a technique that significantly increases the context length of LLMs to an extraordinary 2048k, all while preserving their performance on shorter context spans. Our approach leverages two types of non-uniformities inherent in RoPE positional encoding, identified through an efficient evolutionary search procedure. This strategy yields two major advantages: it supplies optimal initialization for subsequent fine-tuning and also permits an 8$\times$ enlargement of the context window even in the absence of fine-tuning. Expanding on this idea, we put forward a stepwise extension approach—employing LLMs fine-tuned at a 256k context length to ultimately reach a 2048k window, requiring no further fine-tuning. Comprehensive experimental results demonstrate the efficacy of LongRoPE. We anticipate that our LongRoPE-2048k models will unlock a variety of novel long-context applications and stimulate continued advances in the field.

\section*{Broader Impacts} The objective of this research is to contribute to the ongoing progress within the domain of Machine Learning. While our study may have a range of societal implications, we do not identify any that require explicit emphasis in this context.

	\balance

\end{document}